% answers_short.tex —— 2008 年·普通高等学校招生考试·北京卷·文科 · 简短答案速查
% 由 tools/make_short_answers.py 自 answers.tex 的 \gksolhead/\gkanswer 提取，
% 并与 2008/2008_数学真题_北京文_答案要点.md 逐题比对后生成，勿手改（改源 MD 后重跑）。

\begin{gknotice}
\textbf{使用说明}：本节只给最终答案便于快速核对，完整推导见解析版。答案为编写者独立推导，\textbf{非官方标准答案}。
\end{gknotice}

\gkpart{选择题}
\begin{center}\begin{tabular}{ccccccccc}
\toprule
题号 & 1 & 2 & 3 & 4 & 5 & 6 & 7 & 8 \\
\midrule
答案 & D & A & A & C & B & A & C & B \\
\bottomrule
\end{tabular}\end{center}

\gkpart{填空题}
\begin{itemize}[leftmargin=2.6em]
  \item[\textbf{9.}] $\frac43$。
  \item[\textbf{10.}] $\{x\mid x<-2\}$。
  \item[\textbf{11.}] $-8$。
  \item[\textbf{12.}] $10$；$32$。
  \item[\textbf{13.}] $2$；$-2$。
  \item[\textbf{14.}] ②。
\end{itemize}

\gkpart{解答题}
\begin{itemize}[leftmargin=2.6em]
  \item[\textbf{15.}] （1）$\omega=1$；（2）$\left[0,\frac32\right]$。
  \item[\textbf{16.}] （1）证明见解析；（2）$\arccos\frac{\sqrt3}{3}$（约 $54.7^\circ$）。
  \item[\textbf{17.}] （1）$a=0$，$c=2$；（2）$b>0$ 时增区间为 $\mathbf R$；$b<0$ 时增区间为 $(-\infty,-\sqrt{-b})$、$(\sqrt{-b},+\infty)$，减区间为 $(-\sqrt{-b},\sqrt{-b})$。
  \item[\textbf{18.}] （1）$\frac1{40}$；（2）$\frac9{10}$。
  \item[\textbf{19.}] （1）$|AB|=2\sqrt2$，面积 $2$；（2）$y=x-1$。
  \item[\textbf{20.}] （1）$\lambda=3$，$a_1=1$（题面字面；同题版本 $a_3=-3$）；（2）不可能为等差数列；（3）$\lambda\in\bigcup_{k\in\mathbf N^*}(4k^2-2k,\ 4k^2+2k)$。
\end{itemize}
